\section{Two explicit amplification families}\label{sec:amplification}

Let $N_\mu(T)$ count positive primitive triples $A+B=C\le T$ with
$\rad(ABC)<C^\mu$, where $0<\mu<1$. The current estimates
\cite[Proposition~1.1 and Theorems~1.2--1.3]{BBLT} give
\begin{equation}\label{eq:exceptional-exponent}
 N_\mu(T)\ll_{\mu,\delta}T^{F(\mu)+\delta},\qquad
 F(\mu)=\min\left\{\frac{2\mu}{3},\frac{23\mu+3}{40},\frac35\right\},
 \quad\delta>0.
\end{equation}
We test two actual constructions from a primitive seed $a+b=c$.
An amplifier contradicting \eqref{eq:exceptional-exponent} at output
height $T=c^K$, for fixed $K>0$, would need more than
$c^{KF(\mu)+o(1)}$ distinct qualifying outputs. Our conclusions concern
explicit radical certificates, rather than the entire exceptional set.

\subsection{Inherited congruence moduli}
Fix pairwise coprime positive integers $U,V,W$, put $M=UVW$, and let
$R_M=\rad(M)$. Outputs with $U\mid A$, $V\mid B$, $W\mid C$ satisfy
\[
 \rad(ABC)\le\frac{R_MABC}{M}.
\]
Call an output \emph{size-certified} if the right side is at most $C^\mu$.

\begin{proposition}\label{prop:crt-fibre}
For $T\ge1$, there are at most $1+\lfloor\log_2T\rfloor$
size-certified primitive outputs with $C\le T$, for $0<\mu\le1$.
Allowing every inherited template $U\mid a$, $V\mid b$, $W\mid c$
gives at most
\[
 \tau(abc)\bigl(1+\lfloor\log_2T\rfloor\bigr)
\]
distinct outputs. For fixed $K$ and $T\le c^K$ this is $c^{o(1)}$.
\end{proposition}
\begin{proof}
Suppose first $M\ge2$. Certification implies $R_MAB\le M$, and
$R_M\ge2$. Consider two outputs in the strip $L\le A,A'<2L$.
Their determinant $\Delta=AB'-A'B$ is divisible by $UV$ and also by $W$,
since $\Delta=AC'-A'C$. Thus $M\mid\Delta$.
The bounds $B'\le M/(R_MA')$ and $B\le M/(R_MA)$ give
\[
 0<AB',A'B<2M/R_M\le M.
\]
Consequently $|\Delta|<M$, so $\Delta=0$. Coprimality implies
$A\mid A'$ and $A'\mid A$, hence equality of all coordinates.
There is at most one output in each dyadic strip. The strict strip endpoint
also handles $R_M=2$. If $M=1$, certification forces $AB\le1$; only
$(1,1,2)$ can remain, and only when the exponent permits it.

The number of inherited templates is
$\tau(a)\tau(b)\tau(c)=\tau(abc)$. A union bound proves the second claim.
For every fixed $\eta>0$, the elementary estimate
$\tau(n)\ll_\eta n^\eta$ follows from $e+1\le2^e\le p^{\eta e}$
for sufficiently large primes $p$ and bounded multiplicative costs at the
finitely many smaller primes. Since $abc\le c^3$, this proves the last claim.
\end{proof}

Outputs whose remaining factors have additional powers or smaller actual
radicals are not ruled out by this proposition. A count of raw CRT
solutions cannot replace the count of distinct certified primitive triples.

\subsection{Square completion on a conic}
Put $P=abc$ and $R_P=\rad(P)$. Positive integer points on
\begin{equation}\label{eq:square-conic}
 ax^2+by^2=cz^2
\end{equation}
give candidate outputs $(ax^2,by^2,cz^2)$. Retain only primitive outputs.
Then $\gcd(x,y,z)=1$, and the positive coordinates are recovered uniquely
from the output. Their radical is at most $R_Pxyz$.

\begin{lemma}\label{lem:conic-parameters}
For coprime integers $m,n$, not both zero, put
\[
 Q=an^2+bm^2,\quad J=an+bm,\quad
 X=Q-2nJ,\quad Y=Q-2mJ,
\]
and set
\[
 D_0=\gcd(a,m)\gcd(b,n)\gcd(c,m-n),\qquad G=\gcd(X,Y,Q)>0.
\]
Then
\[
 aX^2+bY^2=cQ^2,\qquad
 \gcd(Q,J)=D_0,\qquad G=\gcd(Q,2J),\qquad D_0\mid G\mid2D_0.
\]
Every positive primitive point of \eqref{eq:square-conic} is represented
by $(X/G,Y/G,Q/G)$ for some such pair. Passing to absolute values only
overcounts outputs, and $(m,n)$ and $(-m,-n)$ give the same point.
\end{lemma}
\begin{proof}
Expansion proves the conic identity. Coprimality of $m,n$ gives
$\gcd(Q,2nJ,2mJ)=\gcd(Q,2J)$.
Here is a valuation check of the remaining gcd, including the prime two.
If a prime divides both $Q,J$, the identities
\[
 aQ=J^2-2bmJ+bc m^2,\qquad bQ=J^2-2anJ+ac n^2
\]
show that it divides $abc$. At a prime dividing $a$, the common valuation
is $\min(v_p(a),v_p(m))$: compare the terms of $J=an+bm$ and
$Q=an^2+bm^2$, using that $b,n$ are units whenever $p\mid m$.
The case $p\mid b$ is identical with the roles exchanged.
At $p\mid c$, set $e=v_p(c)$, $r=v_p(n-m)$. If $r=0$ the common
valuation is zero. If $r>0$, both coordinates are units and
$J=a(n-m)+cm$ has valuation $\min(e,r)$ unless $e=r$.
In that exceptional case either $v_p(J)=e$, or
\[
 Q=2mJ-cm^2+a(n-m)^2
\]
has valuation exactly $e$. Thus the common valuation is always
$\min(e,r)$. This argument does not divide by two. It proves the gcd
formula and hence the assertions about $G$.

In the chart $z=1$, a line through $(1,1)$ has equation
$(x,y)=(1,1)+t(n,m)$. Its intersection parameters satisfy
$t(2J+tQ)=0$, giving the displayed second point. Every other rational
point determines such a rational direction, which may be chosen primitive.
The base point itself comes from the tangent direction $(m,n)=(-a,b)$.
Dividing homogeneous coordinates by their gcd gives every primitive point.
\end{proof}

\begin{theorem}\label{thm:conic-fibre}
For every $T>0$, the number of positive primitive square-completion outputs
with $cz^2\le T$ is at most
\begin{equation}\label{eq:conic-fibre}
 \tau(P)\left(1+4\pi\sqrt{T/P}\right).
\end{equation}
The subfamily certified by $R_Pxyz\le(cz^2)^\mu$, for $0<\mu<1$,
has size at most
\begin{equation}\label{eq:conic-certified}
 \tau(P)\left(1+\frac{4\pi}{\sqrt c}T^{\mu/2}\right).
\end{equation}
\end{theorem}
\begin{proof}
Fix the three divisors in $D_0$ from Lemma~\ref{lem:conic-parameters}.
The conditions
\[
 d_a\mid m,\qquad d_b\mid n,\qquad d_c\mid m-n
\]
define a lattice in $\Z^2$ of index $D_0=d_ad_bd_c$. Indeed the first
two have index $d_ad_b$, and the third residue map is surjective there
because these three divisors are pairwise coprime.
The height bound and $G\le2D_0$ imply
\[
 an^2+bm^2=Q\le2D_0\sqrt{T/c}.
\]
This ellipse has area $V=2\pi D_0\sqrt{T/P}$.

A centred ellipse of area $V$ contains at most $1+2V/D_0$ primitive
integer directions in a lattice of index $D_0$, modulo sign. To see this,
select representatives in a fixed half-plane and order them by angle.
There is at most one primitive vector on each positive ray. Consecutive
distinct directions have nonzero determinant divisible by $D_0$, so
their triangles with the origin have area at least $D_0/2$.
These triangles have disjoint interiors and lie in the ellipse. The cases
of zero or one direction satisfy the bound as well. There are $\tau(P)$
divisor templates, proving \eqref{eq:conic-fibre}.

For a certified output write $C=cz^2$. Equation~\eqref{eq:square-conic}
implies $\max(x,y)\ge z$, so $xyz\ge z^2$. Hence
\[
 R_PC/c\le R_Pxyz\le C^\mu,\qquad C^{1-\mu}\le c/R_P.
\]
All certified heights are at most
$U=\min\{T,(c/R_P)^{1/(1-\mu)}\}$. Moreover,
$ab\ge c-1\ge c/2$ and $R_P\ge2$, giving $R_PP\ge c^2$. Therefore
\[
 \sqrt{U/P}
 =U^{\mu/2}\sqrt{U^{1-\mu}/P}
 \le T^{\mu/2}\sqrt{c/(R_PP)}
 \le c^{-1/2}T^{\mu/2}.
\]
Apply \eqref{eq:conic-fibre} at height $U$ to obtain
\eqref{eq:conic-certified}.
\end{proof}

For fixed $K>0$ and $T=c^K$, \eqref{eq:conic-certified} is
$c^{\max\{0,K\mu/2-1/2\}+o(1)}$. Every term defining $F(\mu)$ is
strictly greater than $\mu/2$, so this family cannot exceed the exponent
required to contradict \eqref{eq:exceptional-exponent}.
The conclusion is limited to the certificate. For example, the seed
$(1,8,9)$ and point $(7,2,3)$ give $(49,32,81)$, whose radical is $42$,
while $R_Pxyz=252$. Since $42^2<3^7$,
$42<81^{9/10}$, although the certificate fails even with exponent one.
This example proves that actual exceptional descendants cannot be
identified with the certified subfamily; it is not a disproof of abc.
